% V10 successor manuscript: supplementary M2 argument.
% This appendix is intentionally separate from the frozen author proof.
% Validation status: mathematical argument under Lean review; do not mark
% the concrete M2 endpoint GREEN until its transitive axiom audit passes.
% To include, place \input{appendices/v10-m2-nonedge} in the appendix area.

\section{The non-edge proof of (M2)}\label{app:m2-nonedge}

We give the one-level factorisation argument in the partial-type
representation of Section~6.3.  The only property of the forbidden
family used in the construction is that its non-singleton members are
irreducible; as usual, the empty unary and diagonal record is assumed
not to be forbidden.  In particular, adjoining a vertex carrying this
record and having no $L$-relations to the old vertices does not create
an induced forbidden copy.  The auxiliary relation $E$ is used to
record free levels; it does not belong to the forbidden language $L$.

Write $L^+=L\cup\{E\}$.  A type of level $k$ has ordinary vertices
$0,\ldots,k-1$ and one distinguished type vertex.  Two types of the
same level will be called \emph{compatible} if their restrictions to
ordinary vertices are identical as complete $L^+$-structures.  We
keep both positive and negative atoms throughout.  For a partial
structure $A$, let $\operatorname{tp}^k_A(v)$ denote the complete
$L^+$-type of $v$ over its first $k$ ordinary vertices, provided
$k\leq\operatorname{fl}_A(v)$.

\begin{lemma}\label{lem:m2-retained-record}
Let $F\in M$ and put $\lambda=\widetilde F$.  For every admissible
type $T$ of level $k$, the restriction of $F(T)$ to ordinary
coordinates $\lambda(i)$, $i<k$, and the distinguished type vertex,
transported back along $i\mapsto\lambda(i)$, is exactly $T$ as a
complete $L^+$-type.
\end{lemma}
\begin{proof}
The distinguished singleton record is fixed by the root-extension
property.  The complete record between the type vertex and ordinary
coordinate $i$ is the terminal letter of the canonical successor
which introduces $i$.  This letter occurs unchanged in the weak
image successor, which is an initial segment of $F(T)$.

For ordinary coordinates $i<j<k$, use the canonical successor
introducing $j$ and let $f_j$ be the free cut of its parameter.  If
$i<f_j$, the pair record occurs in the actual parameter, so the weak
successor law transports it to $\lambda(i),\lambda(j)$.  If
$i\geq f_j$, the no-new-tuples clause forces both directed
$L$-relations and both $E$-relations across this pair to be absent.
Strict monotonicity of $\lambda$ puts $\lambda(i)$ beyond the mapped
parameter cut, so these crossing relations are also absent in the
image.  The singleton and diagonal records have already been
transported.  This proves equality of all atoms, including the
negative ones.
\end{proof}

\begin{lemma}\label{lem:m2-simultaneous-image}
Let $A$ be a finite admissible partial structure and let
$v_1,\ldots,v_r$ be ordinary vertices whose free levels are at least
$n>0$.  Put $T_j=\operatorname{tp}^n_A(v_j)$.  There exists an
admissible type $X$ with $A$ in its ordinary socle such that, for
every $j$,
\[
  \operatorname{tp}^{\lambda(n)}_{F(X)}\bigl(\lambda(v_j)\bigr)
       =F(T_j).
\]
Moreover, all $L^+$-relations among the retained ordinary vertices
of $A$ occur unchanged among the corresponding coordinates of $F(X)$.
\end{lemma}
\begin{proof}
Extend $A$ by a neutral ordinary buffer vertex and then by a
neutral distinguished vertex, using an incoming $E$-column whose free
cut is beyond the whole ordinary socle.  The added vertices have no
$L$-relations to old vertices, and may be spaced so that the $E$
axioms hold.  Irreducibility, including the allowed-neutral-singleton
condition, shows that the extension is admissible.  Let $X$ be its
complete type.

For $v_j$, let $f_j\geq n$ be its actual free cut.  As $f_j>0$, the
step introducing $v_j$ has the nonempty canonical parameter
$P_j=\operatorname{tp}^{f_j}_A(v_j)$, and
$T_j=P_j\restriction n$.  The weak successor law for $F$ puts the
complete image parameter $F(P_j)$ at ordinary coordinate
$\lambda(v_j)$, below the image of the entire containing type $X$.
Its restriction through $\lambda(n)$ is $F(T_j)$ by prefix
preservation.  This proves the displayed identity simultaneously for
all $j$.  The remaining assertion is
Lemma~\ref{lem:m2-retained-record} applied to $X$.
\end{proof}

\begin{lemma}\label{lem:m2-compatibility}
If $T$ and $U$ are compatible admissible types of the same positive
level, then $F(T)$ and $F(U)$ are compatible.
\end{lemma}
\begin{proof}
Realise their common ordinary $L^+$-socle once, and add two vertices
realising respectively the distinguished-vertex types of $T$ and
$U$.  Put no $L$-relations between these two new vertices.  If
necessary insert neutral padding between them and the preceding
vertices to obtain valid $E$-columns and free cuts at least the
common level.  Any nontrivial irreducible forbidden copy uses at most
one of the two new vertices and hence is already excluded by the
admissibility of $T$ or $U$; a forbidden singleton is excluded as
well.  By Lemma~\ref{lem:m2-simultaneous-image}, $F(T)$ and $F(U)$
are both the complete types of ordinary vertices above one and the
same initial $L^+$-socle in $F(X)$.
\end{proof}


\paragraph{Bounded upper witnesses.}
The finite-age verification below can be shortened by considering a
single induced forbidden copy rather than every upper vertex of an
arbitrary age witness.  Indeed, suppose the conclusion fails and
choose an induced copy $e$ of some $H\in\mathcal F$ in $W$.
Since the initial socle through $m$ is induced from an admissible
target and $W\setminus\{m\}$ is $\mathcal F$-free, the copy contains
$m$ and at least one vertex above $m$.  It uses at most $|H|$
upper vertices.

Keep the initial target socle and only those upper vertices $e(a)>m$
in the age calculation.  The other numerical upper positions can be
filled by vertices carrying the allowed neutral singleton and no
incident $L$-relation.  This does not create an irreducible forbidden
configuration: a forbidden copy containing a neutral filler would
violate irreducibility, and a forbidden singleton cannot be neutral.
Consequently one may use exactly $n+|H|$ original source positions:
the first $n$ are mapped to retained coordinates
$\lambda(0),\ldots,\lambda(n-1)$, while the position $n+a$
is retained only when $e(a)>m$, in which case it represents $e(a)$.
The map is increasing on these genuine positions, since $e$ and
$\lambda$ are increasing and all upper images exceed $m$.
The corresponding source $L$-structure is an induced pullback of
$W\setminus\{m\}$ on genuine positions, with neutral non-edges
elsewhere, and is therefore $\mathcal F$-free.

This bounded construction preserves every required upper--upper
$L$-atom, including absent relations.  It does not request an
$E$-expansion of $W$.  Only after the source $L$-structure is obtained
do we apply the non-edge padding construction, with the one common
ordinary $E$-socle of the prescribed source types.

\begin{lemma}\label{lem:m2-compatible-age}
Let $F\in M$, let $n>0$, and suppose that
$\lambda(n-1)<m=\lambda(n)-1$.  Define $D$ by deleting ordinary
coordinate $m$ from $F(T)$ for $\ell(T)\geq n$ and leaving $F(T)$
unchanged below level $n$.  Put
$\mathcal S=D[\mathcal T(n)]$ and prescribe
$f(D(T))=F(T)$ for $T\in\mathcal T(n)$.

Fix a compatibility class of $\mathcal S$, represented by $S_0$.
Let $W$ be a finite ordered $L$-structure on an initial segment
whose first $m+1$ vertices induce the ordinary $L$-socle of $f(S_0)$.
Suppose that every upper vertex $x>m$ has its complete $L$-type over
$0,\ldots,m$ equal to the $L$-reduct of some $f(S_x)$, where
$S_x\in\mathcal S$ is compatible with $S_0$.  If
$W\setminus\{m\}$ avoids the forbidden family, then so does $W$.
\end{lemma}
\begin{proof}
If $W$ has no upper vertex, it is the ordinary socle of an
admissible type.  Otherwise enumerate the upper vertices of $W$ as
$x_1<\cdots<x_r$ and choose source types $T_j$ with
$S_{x_j}=D(T_j)$.  The same $T_j$ may occur more than once.  By
compatibility and Lemma~\ref{lem:m2-retained-record}, the types
$T_j$ share one complete ordinary $L^+$-socle of level $n$; the
$L$-type of each distinguished vertex over the first $n$ source
coordinates is the pullback of the given type of $x_j$ over the
retained target coordinates $\lambda(i)$, $i<n$.

Form an ordered $L$-structure $A$ whose initial vertices are
$0,\ldots,n-1$ and which has an upper vertex $v_j$ corresponding to
$x_j$.  Copy every $L$-atom, in both directions, from the induced
substructure of $W\setminus\{m\}$ on the coordinates
\[
  \{\lambda(i):i<n\}\cup\{x_1,\ldots,x_r\}.
\]
Thus $A$ avoids the forbidden family, its initial $L$-socle is the
common initial $L$-socle of the $T_j$, and its upper vertices have
the required complete $L$-types over that socle.  In particular,
\emph{all relations between distinct upper vertices are retained}.

Pad $A$ by keeping its first $n$ vertices fixed and moving each
upper vertex $v_j$ to
\[
       a_j=n+1+2(j-1).
\]
Every new position receives the allowed neutral singleton and no
$L$-relations.  The padded $L$-structure still avoids the forbidden
family: a nontrivial irreducible forbidden copy cannot contain a
padding vertex, and its induced copy on the old vertices would
project to $A$.

Equip the first $n$ positions with the common prescribed ordinary
$E$-socle of the $T_j$.  Give each real upper vertex $a_j$ the
maximal spaced incoming $E$-column, namely
\[
       E(u,a_j)\quad\Longleftrightarrow\quad u+1<a_j,
\]
and give each padding vertex an empty incoming column.  This is a
valid partial structure.  In particular the complete type of the
real vertex $a_j$ over the first $n$ vertices is exactly $T_j$:
the $L$-record was copied, the ordinary $E$-socle was fixed, and all
$E$-incidences of the distinguished vertex are forced by the cut
$n$.  Notice that the padding vertices are $L$-isolated, not
necessarily $E$-isolated.

Apply Lemma~\ref{lem:m2-simultaneous-image} to this single padded
structure.  We obtain an admissible type $X$ such that
\[
  \operatorname{tp}^{\lambda(n)}_{F(X)}(\lambda(a_j))=F(T_j)
     \qquad(1\leq j\leq r).
\]
Define a map from the vertices of $W$ to ordinary vertices of $F(X)$
by
\[
  e(z)=\begin{cases}
       z,&z\leq m,\\
       \lambda(a_j),&z=x_j.
       \end{cases}
\]
This map is increasing and injective.  Indeed,
$\lambda(a_j)>\lambda(n)=m+1$ and the $a_j$ are strictly increasing.
We check that it is an induced $L$-embedding.

On the initial $m+1$ vertices, the displayed type identity gives
exactly the ordinary socle of $F(T_1)$, which agrees with that of
$F(T_j)$ for every $j$ by Lemma~\ref{lem:m2-compatibility}.
Between an initial vertex and $x_j$, the same type identity gives
both directed relations and their complements, including the skipped
coordinate $m$.  Between two upper vertices, use
Lemma~\ref{lem:m2-retained-record} for $X$: the relation between
$\lambda(a_i)$ and $\lambda(a_j)$ is precisely the relation between
$a_i$ and $a_j$, and therefore between $x_i$ and $x_j$.  Singleton
and diagonal records are treated identically.  Thus $W$ is an
induced substructure of the ordinary $L$-socle of $F(X)$.  This
socle avoids the forbidden family; hence so does $W$.
\end{proof}

\begin{proof}[Proof of (M2)]
Let $F\in M$, choose the first source level $n$ relevant to the
desired factorisation and put $\lambda=\widetilde F$.  Suppose
$\lambda(n-1)<\lambda(n)-1$ and set $m=\lambda(n)-1$.
Since $\lambda(0)=0$, we have $n>0$ and $m>0$.

First define $D$ by removing coordinate $m$ from the image types
above level $n$, as in Lemma~\ref{lem:m2-compatible-age}.  This is a
shape-preserving map with level map
\[
 \widetilde D(k)=\begin{cases}
       \lambda(k),&k<n,\\
       \lambda(k)-1,&k\geq n.
       \end{cases}
\]
Deletion preserves the forbidden-free $L$-age and commutes with
initial-socle restriction.  At the crossing level, the image of the
source successor has parameter cut strictly below the deleted
coordinate, so removing that coordinate neither creates an
$E$-crossing nor changes its terminal successor letter.  At all
other levels the weak successor law follows from the corresponding
image successor by deletion.  The level map is strictly increasing,
root extensions are unchanged, and injectivity follows by recovering
the retained source coordinates.  Consequently $D\in M$.

At level $n$ the removed coordinate is the last ordinary vertex,
so $D(T)=F(T)\restriction m$.  Since $D$ is injective, the
prescription
\[
       f(D(T))=F(T)\qquad (T\in\mathcal T(n))
\]
is well defined and satisfies the predecessor condition (B1).
It also satisfies (B2): if two deleted images have identical ordinary
$L^+$-socles, Lemma~\ref{lem:m2-retained-record} recovers the
corresponding source socles, and
Lemma~\ref{lem:m2-compatibility} shows that the full target socles
are identical.  The targets are admissible, being genuine images
under $F$.

Lemma~\ref{lem:m2-compatible-age} verifies the finite age test in
every common ordinary $L^+$-compatibility class of prescribed source
types.  This is the compatibility-restricted form of (B3) required
by the boring-extension construction: every upper source type
arising at an insertion stage belongs to one such class because it
is read from one ambient partial structure.  Therefore that
construction supplies a single $G\in M$ which skips precisely $m$
and satisfies
\[
       G(D(T))=F(T)\qquad (\ell(T)=n).
\]
For $\ell(T)<n$, $D(T)=F(T)$ and
$\ell(F(T))\leq\lambda(n-1)<m$.  A shape-preserving map which skips
only $m$ fixes all nodes below level $m$ (by the root-extension and
successor properties), hence $G(D(T))=F(T)$ at these levels as well.
Thus
\[
       (G\circ D)(T)=F(T)\qquad (\ell(T)\leq n).
\]
Taking $F_1=D$ and $F_2=G$ establishes (M2).
\end{proof}

% \todo[inline]{\textbf{Řehořek: M2 validation OPEN.} The source padding,
% parameter transport, partial inverse and conditional global extension
% have independent Lean implementations. The missing end-to-end checks
% are compatibility for the actual boundary data, the finite induced
% embedding in Lemma~\ref{lem:m2-compatible-age}, and the final
% concrete M2 endpoint with a transitive standard-axiom audit.}
